% Notes of the Montreal seminar, 1970 (cote 129, typed part): archive cote 129, site page 65 (PDF page 66).
% typescript. Draft compared: none.
% AI-transcribed by Claude Opus 5.5 (claude-opus-5-5, adaptive thinking, effort high), Message Batches API, batch pilotnext (msgbatch_01A9sdu8SYq23bEtqX9tUncQ), 2026-11-13; not yet checked by eye.
\scanpage{129/65}
\mspage{\hand{III-10}}

\margin{X} correspondre \deleted{G univoquement}\inserted{biunivoquement} le système projectif p-adique
\[ T_p(G) = (G(n)_{n\geq 1} \; ; \; G(n+1) \xrightarrow{p} G(n)) \]
où la multiplication par $p$ identifie $G(n)$ à $G(n+1)/p^nG(n+1)$ et
où $G(1)$ est annulé par $p$ et est représentable par un schéma fini
localement libre sur $S$.

5.5. \uline{Dual d'un groupe de Barsotti-Tate}. A tout groupe de BT
$G$ sur $S$, on associe donc le système p-adique $T_p(G)$. Ce système
projectif définit par dualité de Cartier un système co-p-adique de
groupes finis localement libres $G(n)^*$, annulés par $p^n$ et plats sur
$\Lambda_n$ (2.2). Par définition, le dual de $G$ est le groupe de BT $G^*$ as-%
socié à ce système co-p-adique:
\[ G^* = \lim_{\substack{\rightarrow \\ n}} G(n)^* \; . \]

On a en particulier $G^*(n) = G(n)^*$, et on obtient ainsi un
foncteur contravariant, compatible aux changements de base, qui définit
une équivalence entre la catégorie $BT(S)$ des p-groupes de Barsotti-%
Tate sur $S$ et son opposée, car on a un isomorphisme fonctoriel canonique
\[ G \xrightarrow{\sim} (G^*)^* \; . \]

5.6. \uline{Module de Dieudonné associé à un groupe de BT sur un corps}
\uline{parfait}. Reprenant les notations de II, 4.1 et désignant par $BT(k)$
la catégorie des p-groupes de BT sur le corps parfait $k$, on a le

\uline{Théorème} (cf. II, 5.3.1) Il existe une équivalence de catégories

\mspage{\hand{65}}
